% SECTION_GHOSTS independent audit: fixed-pair compactness and continuous cohomology.
% No assertion in this file settles the ordinary all-depth support question.
\subsection{The fixed-pair compatibility theorem}

Let $G=G_{\mathbb Q}$ and let $F_\ell$ be the free pro-$\ell$
geometric fundamental group of the tripod, based at $\overrightarrow{01}$.
We always use the actual arithmetic action obtained from that tangential
section. Put $N_{r,\ell}=F_\ell/\Gamma_{r+1}F_\ell$ and
$L_{r,\ell}=\Gamma_rF_\ell/\Gamma_{r+1}F_\ell$, with closed
commutator subgroups. The associated graded action is
$L_{r,\ell}=\mathbb Z_\ell(r)^{b_r}$: a weight-$r$ bracket is
multilinear in weight one, on which $G$ acts through $\chi_\ell$.

\begin{result}{A.20}{Proved}{Compactness with the abelian cocycle fixed}
\label{coh:compact}
For $r\geq2$, $Z^1_{\rm cts}(G,\mathbb Z_\ell(r))$ is a finite
free $\mathbb Z_\ell$-module, compact for the compact-open topology.
Fix a literal cocycle $a\in Z^1_{\rm cts}(G,N_{1,\ell})$.
Then the space
\[
 \mathscr S_{r,\ell}(a)=
 \{c\in Z^1_{\rm cts}(G,N_{r,\ell}):c_{\rm ab}=a\}
\]
is compact Hausdorff. More precisely, it is parametrized continuously by
a closed subset of a finite power of $\mathbb Z_\ell$, with Hall
coordinates polynomial over $\mathbb Q_\ell$ in the parameters and
continuous cochains as coefficients. This confirms
GT\_TOWER\_CHECK C.1.3--C.1.6. No compactness is asserted when $a$
varies.
\end{result}
\begin{proof}
We first prove the assertion about additive cocycles. At a prime
$q\ne\ell$, inertia acts trivially on $T=\mathbb Z_\ell(r)$.
Restriction of a cocycle to inertia is therefore a homomorphism. Wild
inertia and the prime-to-$\ell$ tame part have zero image. If $\tau$
is a tame $\ell$-generator and $\sigma$ is arithmetic Frobenius,
then the cocycle identity and $\sigma\tau\sigma^{-1}=\tau^q$ give
$(q^r-q)c(\tau)=0$. Since $r\geq2$ and $T$ is torsion free,
$c(\tau)=0$. The cocycle kills the closed normal subgroup generated by
all inertia groups outside $S=\{\ell,\infty\}$, and descends to
$G_{\mathbb Q,S}$.

For every $n$, the set $Z^1(G_{\mathbb Q,S},T/\ell^n)$ is finite.
Indeed a cocycle defines a homomorphism to the finite semidirect product
$(T/\ell^n)\rtimes\operatorname{im}(\chi_\ell^r\bmod\ell^n)$.
Its fixed field has bounded degree and ramifies only in $S$.
Milne's fixed-ramification finiteness theorem
\cite[Theorem 8.42, p.151]{CohMilneANT} gives finitely many such fields,
and each finite group has finitely many maps into the displayed target.
Compatible finite reductions are exactly continuous $T$-valued cocycles;
thus their inverse limit $Z$ is compact. Reduction has kernel $\ell Z$,
because any cocycle with values in $\ell T$ can be divided by $\ell$.
Hence $Z/\ell Z$ is finite. Lifting a basis and successively correcting
modulo $\ell^n$ expresses every cocycle as a $\mathbb Z_\ell$-linear
combination of finitely many lifts. The residual cocycles tend uniformly
to zero. Consequently $Z$ is finitely generated and torsion free, hence
free. Its module topology agrees with compact-open topology: on a compact
domain, congruence modulo $\ell^n$ is a uniform neighbourhood condition.

We record the coordinate input rather than assume that integral and
rational nilpotent groups coincide. The free discrete nilpotent group of
class $r$ has its integral Hall basis. Collection gives polynomial
multiplication and powering with rational coefficients; denominators are
bounded at fixed $r$ \cite[\S\S3--4]{CohCantEick2019}. Extending these
formulas by continuity to $\mathbb Z_\ell^{h_r}$ gives a compact
nilpotent pro-$\ell$ group. Its Hall one-parameter subgroups are images
of $\mathbb Z_\ell$. Every finite quotient is nilpotent and generated
by elements of $\ell$-power order, and is therefore an $\ell$-group.
Conversely, a homomorphism of the discrete group to a finite $\ell$-group
of exponent $\ell^e$ kills the coordinate box with all Hall exponents
in $\ell^e\mathbb Z$. It is continuous for this coordinate topology
and extends uniquely to the compact group. This universal property
identifies the group with $N_{r,\ell}$, including $\ell=2,3$.
Deleting the last-weight coordinates is the quotient map; setting them
equal to zero is a continuous section as a map of spaces. Substitution
of the actual images of the two generators gives action coordinates
polynomial in the input coordinates, with continuous coefficients in $g$.

The remaining point is closedness of the locus of soluble defects.
Given finitely many continuous cochains
$c_i\in C^2(G,T\otimes\mathbb Q_\ell)$, put
\[
 R=\{u\in\mathbb Z_\ell^m:\textstyle\sum_i u_ic_i
                    \in\delta C^1(G,T)\}.
\]
It is a submodule of a finite free module over the Noetherian PID
$\mathbb Z_\ell$, hence is free on finitely many vectors $r_j$.
The corresponding injective map $\mathbb Z_\ell^k\to
\mathbb Z_\ell^m$ is continuous; its compact image is $R$, so $R$
is closed. Choose a continuous integral primitive $h_j$ for the
combination determined by each $r_j$. A rational left inverse of the
basis matrix supplies linear forms $\lambda_j$, integral on $R$.
Then $h(u)=\sum_j\lambda_j(u)h_j$ is a continuous choice of primitive
on $R$. Notice that we did not assert that all global coboundaries form
a closed subspace of the unrestricted continuous cochain space.

Now induct on $r$. The initial parameter space is the one point $a$.
For a polynomially parametrized family $c_t$ at class $r$, lift through
the zero-last-coordinate Hall section. Its central defect is a finite
sum $\sum_i v_i(t)c_i$, with $v_i$ integral monomials. It is soluble
exactly on the closed locus $v(t)\in R$. Subtracting $h(v(t))$ gives
one lift there. All other lifts differ by a literal cocycle in
$Z^1(G,L_{r+1,\ell})$, a finite free compact module by the first
paragraph. Adjoining those parameters completes the induction. The
resulting evaluation map of the compact parameter space times $G$ is
continuous; its adjoint into the compact-open function space is
continuous and surjective onto $\mathscr S_{r,\ell}(a)$.
\end{proof}

\begin{result}{A.21}{Proved}{The exact limit equivalence and its quantifiers}
\label{coh:limit}
Fix $(A,B)\in(\widehat{\mathbb Q^{\times}})^2$. For the tripod over
$\mathbb Q$, the following are equivalent:
\begin{enumerate}
\item the pair is the abelianization of a continuous section of
      $E_U^{\rm nil}\to G_{\mathbb Q}$;
\item for every prime $\ell$, its $\ell$-component is the
      abelianization of a continuous section of the full geometrically
      pro-$\ell$ quotient;
\item for every positive integer $r$, the same pair is the
      abelianization of a continuous section of the integral class-$r$
      quotient;
\item for every $\ell$ and every $r$, the same $\ell$-component
      survives the pro-$\ell$ class-$r$ quotient.
\end{enumerate}
The corresponding fixed-$\ell$ equivalence holds. It continues to hold
after imposing geometric conjugacy to one prescribed compatible family
of local point sections, or local cuspidal sections, at every level.
This confirms GT\_TOWER\_CHECK C.1.7. The assertion does not permit
the pair to change with $r$ or $\ell$.
\end{result}
\begin{proof}
Every finite $\ell$-group is nilpotent. Thus
$F_\ell\xrightarrow{\sim}\varprojlim_rN_{r,\ell}$ as profinite
groups. The maximal pronilpotent quotient of the geometric free
profinite group is $\prod_\ell F_\ell$: the finite quotients on
both sides are precisely the finite nilpotent quotients, decomposed
into their Sylow factors. The arithmetic action preserves these
characteristic factors. Similarly the integral class-$r$ quotient is
$\prod_\ell N_{r,\ell}$, since its finite quotients are precisely
the finite nilpotent groups of class at most $r$. These are continuous,
surjective, $G$-equivariant quotient maps.

Choose a literal abelian cocycle representing the prescribed component.
A section with that abelian cohomology class can be geometrically
conjugated to have the prescribed literal cocycle: the discrepancy is a
coboundary, and its defining abelian element lifts through
$N_{r,\ell}\twoheadrightarrow N_{1,\ell}$. If every class is
soluble, the spaces of normalized cocycles in A.20 are nonempty compact
spaces. In their product, adjacent-level compatibility is a closed
condition. Every finite collection of these conditions is soluble by
choosing a section at the largest level and projecting. Compactness
therefore gives a simultaneous compatible sequence. Its map into the
inverse limit is continuous, since each component map is continuous.
Taking the product over $\ell$ gives a continuous cocycle into the
pronilpotent group. Projections give all converse implications.

For a prescribed local section, its geometric-conjugacy orbit at a
fixed class is the continuous image of the compact geometric group
in the local compact-open function space. The orbit is closed.
Restriction of global cocycles is continuous, so any collection of
local matching conditions defines a closed subspace of A.20. The same
compactness proof applies. At each place, nonempty compact sets of
conjugators also have the finite-intersection property; their limit is
an actual conjugator at full pro-$\ell$ level. Products over $\ell$
give a pronilpotent conjugator. Conjugacy throughout is geometric;
we never quotient by the full arithmetic group.
\end{proof}

For ordinary pairs, each valuation is an actual integer. Hence a nonzero
support valuation remains nonzero in every $\mathbb Z_\ell$.
An affirmative finite-support theorem for any single coefficient prime
would therefore imply the full pronilpotent finite-support theorem.
The converse is not a logical consequence: the full theorem concerns
the intersection of the pro-$\ell$ survivor sets, not any one of them.
This last observation specifies the logical relation only; it proves no
finite-support assertion.

\subsection{The continuous cohomology that a lifting argument must retain}

\begin{result}{A.22}{Proved}{Strict integral localization}
\label{coh:strict}
Let $r\geq2$, let $T=\mathbb Z_\ell(r)$, and put
$S=\{\ell,\infty\}$. After compatible Tate coordinates are chosen,
define
\[
 R_\ell=\{(z_q)_{q\ne\ell}\in\prod_{q\ne\ell}\mathbb Z_\ell:
                         z_q\longrightarrow0\},
 \qquad (d_rz)_q=(q^{r-1}-1)z_q.
\]
There is a natural exact sequence of abelian groups
\[
 0\longrightarrow H^2_{\rm cts}(G_{\mathbb Q,S},T)
 \longrightarrow H^2_{\rm cts}(G,T)
 \xrightarrow{\Psi_{\ell,r}}R_\ell/d_rR_\ell
 \longrightarrow0.
\]
The denominator is the actual image, without closure. At $r=3$, the
first group is zero for every $\ell$, so $\Psi_{\ell,3}$ is an
isomorphism. This confirms GT\_TOWER\_CHECK C.1.13.
\end{result}
\begin{proof}
Write $T_n=T/\ell^n$. The localization triangle for removing finitely
many closed points from $X=\operatorname{Spec}\mathbb Z[1/\ell]$,
followed by the filtered union over those sets, is
\[
 R\Gamma(G_{\mathbb Q,S},T_n)\longrightarrow
 R\Gamma(G,T_n)\longrightarrow
 \bigoplus_{q\ne\ell}R\Gamma(\mathbb F_q,T_n(-1))[-1]
 \longrightarrow R\Gamma(G_{\mathbb Q,S},T_n)[1].
\]
The maps are the coefficient-compatible supports morphisms. The
purity input is $R^2i^!\mu_{\ell^n}=\mathbb Z/\ell^n$, with
the other derived sheaves zero, tensored by the remaining twist;
the exact primary locators are
\cite[II.1.7(b), pp.151--152; II.2.9, pp.170--171]{CohMilneADT}.
Every finite-valued continuous Galois cochain factors through a finite
Galois quotient, hence occurs after removing finitely many primes.
This identifies the middle term with the absolute continuous cochain
complex used here.

At each residue field the coefficient-compatible procyclic resolution
is $[\mathbb Z/\ell^n\xrightarrow{q^{r-1}-1}
\mathbb Z/\ell^n]$ in degrees zero and one. If the prime-to-$\ell$
part of Frobenius acts nontrivially, the scalar is a unit and this
complex is contractible; otherwise it is the usual pro-$\ell$ cyclic
resolution. Thus the last term before the shift is represented by
\[
 [\,\bigoplus_{q\ne\ell}\mathbb Z/\ell^n
                   \xrightarrow{d_r}
       \bigoplus_{q\ne\ell}\mathbb Z/\ell^n\,].
\]
Its transition maps are degreewise surjective and its inverse limit is
$[R_\ell\xrightarrow{d_r}R_\ell]$. A compatible sequence belongs
to this restricted module exactly because its reduction at every
coefficient exponent has finite support.

For completeness, the derived inverse limit is the mapping fibre of
one minus shift on the product of the complexes. This construction
preserves the displayed triangles. Global and fixed-ramification
continuous cochain complexes also reduce surjectively in every degree,
by choosing lifts of their finitely many values. Their derived limits
are therefore the ordinary continuous $T$-cochain complexes; this is
also Jannsen's comparison
\cite[(2.1), Theorem (2.2), p.215]{CohJannsen1988}.
The relevant exact sequence is
\[
 \ker d_r\longrightarrow H^2(G_{\mathbb Q,S},T)
 \longrightarrow H^2(G,T)\longrightarrow\operatorname{coker}d_r
 \longrightarrow H^3(G_{\mathbb Q,S},T)
 \longrightarrow H^3(G,T).
\]
The first term is zero by torsion freeness. The last arrow is injective:
finite coefficient $H^2$ at fixed ramification is finite, so its
derived limit vanishes; finite coefficient degree-three cohomology
is identified with real cohomology by
\cite[I.4.10(c), p.57]{CohMilneADT}. The same holds in the integral
limit because the real degree-two groups are finite. Real restriction
therefore identifies $H^3(G_{\mathbb Q,S},T)$ with its real group,
and factors through $H^3(G,T)$. This proves the short exact sequence,
including $\ell=2$.

Finally $K_4(\mathbb Z)=0$. Localizing at $\ell$ does not change
its $\ell$-primary part: the localization sequence has residue group
$K_3(\mathbb F_\ell)$ of order $\ell^2-1$, prime to $\ell$
\cite[IV.1.13; V.6.8]{CohWeibelKbook}. The odd-primary comparison and the
two-primary comparison, respectively
\cite[VI.8.2; VI.9.11, case $n=4$; Table VI.10.1.1]{CohWeibelKbook},
give $H^2(G_{\mathbb Q,S},\mathbb Z_\ell(3))=0$.
In the two-primary formula the signature parameter satisfies
$0\leq\rho<r_1=1$, hence is zero. This proves the final assertion.
\end{proof}

\begin{result}{A.23}{Proved}{A nonzero class invisible at every finite absolute coefficient level}
\label{coh:691}
For $T=\mathbb Z_{691}(12)$, inflation identifies
$H^2(G_{\mathbb Q,\{691,\infty\}},T)\simeq\mathbb Z/691$
with a nonzero subgroup of $H^2_{\rm cts}(G,T)$. Every element of
this subgroup restricts to zero at every place and reduces to zero in
every $H^2(G,T/691^n)$. A nonzero element defines a continuous
profinite extension by $T$ which splits locally and modulo every
$691^n$, but does not split continuously over $G$.
This confirms GT\_TOWER\_CHECK C.1.8--C.1.9. The extension is not
asserted to be a quotient of the tripod fundamental group.
\end{result}
\begin{proof}
The calculation $K_{22}(\mathbb Z)=\mathbb Z/691$ and the odd-primary
comparison in \cite[VI.8.2, VI.10.3]{CohWeibelKbook} give the
fixed-ramification group. A.22 proves injection of inflation.
At $q\ne691$ an inflated class factors through the unramified
procyclic quotient, whose integral second cohomology is zero: the
finite second cohomology vanishes and finite first cohomology has no
derived limit. At $691$, local duality identifies the dual of the
integral second group with
$(\mathbb Q_{691}/\mathbb Z_{691})(-11)^{G_{\mathbb Q_{691}}}$.
A Teichm\"uller unit acts on this group by a scalar different from
one modulo $691$, since $690\nmid11$, so the invariants vanish.
The real group vanishes by averaging with odd-primary coefficients.

For finite $T_n=T/691^n$, its Tate dual is
$\mathbb Z/691^n(-11)$ and has cyclic splitting field. A globally
defined first cohomology class in that dual module which vanishes
locally everywhere is zero: after restriction to the splitting field
it is a homomorphism to a finite constant module, killed by
Chebotarev; it consequently inflates from the finite cyclic splitting
group. A prime whose Frobenius generates that group then detects the
inflated class, because inflation on $H^1$ is injective. The dual
localization kernel is therefore zero. Poitou--Tate duality
\cite[I.4.10(a), pp.56--57]{CohMilneADT} makes localization on
$H^2(G,T_n)$ injective. All finite reductions of our integral class
have zero localizations, and hence are zero.

Choose a normalized continuous cocycle $c$ representing a nonzero
class. On $T\times G$ use the multiplication
$(u,g)(v,h)=(u+gv+c(g,h),gh)$. It is a compact totally disconnected
group; the cocycle identity gives associativity. A continuous section
has the form $g\mapsto(b(g),g)$ and exists exactly when
$\delta b=-c$. Nonvanishing rules this out globally, whereas the
proved local and finite-coefficient vanishings provide the asserted
splittings. Their choices need not be compatible as the coefficient
exponent varies.
\end{proof}

\begin{result}{A.24}{Proved}{The derived-limit kernel and an explicit restricted-product defect}
\label{coh:limone}
For odd $\ell$ and $r\geq2$, the kernel of integral localization
on $H^2(G,\mathbb Z_\ell(r))$ is
$\varprojlim{}^1_nH^1(G,\mathbb Z/\ell^n(r))$.
For $(\ell,r)=(3,2)$ it contains a subgroup isomorphic to
\[
       \frac{\prod_{j\geq1}\mathbb Z_3}{c_0(\mathbb Z_3)},
 \qquad c_0(\mathbb Z_3)=\{(x_j):x_j\to0\}.
\]
This confirms GT\_TOWER\_CHECK C.1.10. The displayed subgroup is
uncountable; it is not identified with the obstruction of any ordinary
pair.
\end{result}
\begin{proof}
Jannsen's exact sequence, or directly one minus shift on the continuous
cochain complexes, is
\[
 0\longrightarrow\varprojlim{}^1H^1(G,T_n)
 \longrightarrow H^2_{\rm cts}(G,T)
 \longrightarrow\varprojlim H^2(G,T_n)\longrightarrow0.
\]
For odd $\ell$, the splitting group of
$T_n^D=\mathbb Z/\ell^n(1-r)$ is cyclic. The Chebotarev and duality
argument of A.23 proves finite-coefficient localization injective.
At each completion the finite first cohomology groups are finite, so
their derived limit is zero. An integral class is therefore locally
zero exactly when every finite absolute reduction is zero. This proves
the first assertion, with the coefficient transition maps retained.

At $(3,2)$, the fixed-ramification integral second group is zero by
$K_2(\mathbb Z)\{3\}=0$, and degree three is zero. Its finite
second groups are consequently zero. The finite residue sequence is
\[
 0\to H^1(G_{\mathbb Q,\{3,\infty\}},T_n)
 \to H^1(G,T_n)\to
 D_n:=\bigoplus_{q\ne3}H^0(\mathbb F_q,\mathbb Z/3^n(1))\to0.
\]
It follows from the same purity triangle as A.22. The left groups are
finite, and thus $\varprojlim{}^1H^1(G,T_n)=\varprojlim{}^1D_n$.

Choose distinct primes $q_j\equiv1+3^j\pmod{3^{j+1}}$.
Chebotarev for a cyclotomic field supplies them. Put $d_j=q_j-1$;
then $v_3(d_j)=j$. Projection to those primes splits the inverse
subsystem
$D'_n=\bigoplus_j\ker(d_j:\mathbb Z/3^n\to\mathbb Z/3^n)$
of $D_n$. Consider the two-term complexes
$P_n=[\bigoplus_j\mathbb Z/3^n\xrightarrow{d}
\bigoplus_j\mathbb Z/3^n]$ in degrees zero and one. Their inverse
limit is $[R\xrightarrow{d}R]$, where $R=c_0(\mathbb Z_3)$,
and the derived-limit exact sequence identifies
$\varprojlim{}^1D'_n$ with the kernel of
$R/dR\to\varprojlim H^1(P_n)$.

The kernel consists exactly of classes of $z\in R$ such that
$z_j\in d_j\mathbb Z_3$ for all $j$. In one direction this is
immediate. In the other, writing $z_j=d_jy_j$, a primitive modulo
$3^n$ may be taken with only $j<n$ nonzero; hence all finite
cohomology classes vanish. Because $v_3(d_j)\to\infty$, multiplication
by $d$ embeds the entire product $\prod_j\mathbb Z_3$ in $R$.
Its inverse image of $dR$ is exactly $R$. This proves the asserted
quotient identification. Binary sequences define the same class only
when they differ at finitely many coordinates, so there are continuum
many classes. The sequence $z_j=d_j$ is a particular nonzero class:
its only coordinatewise primitive is the constant sequence $1$, which
does not tend to zero. The script \path{code/check_cohomology_audit.py}
checks finite instances of this description; the infinite assertion
follows from the argument just given.
\end{proof}

\subsection{The cup obstruction and the twenty-four unrestricted lifts}

\begin{result}{A.25}{Proved}{Ordinary pairs have zero continuous cup and a class-two section}
\label{coh:cup}
Every ordinary completed pair $(A,B)$ satisfies
$A\cup B=0$ in
$H^2_{\rm cts}(G_{\mathbb Q},\widehat{\mathbb Z}(2))$.
It is therefore the abelianization of a continuous section of the actual
integral class-two quotient of the tripod. This confirms
GT\_TOWER\_CHECK C.1.11--C.1.12. It does not assert local geometric
matching of that section.
\end{result}
\begin{proof}
Fix $\ell$ and literal Kummer cocycles $a,b$ for the two components.
First their cup is zero at every finite coefficient exponent. Indeed
the norm-residue isomorphism identifies finite weight-two Galois
cohomology with $K_2(\mathbb Q)/m$. The tame-symbol exact sequence
has kernel $K_2(\mathbb Z)=\mathbb Z/2$, generated by $\{-1,-1\}$
\cite[III.5.2.2, III.6.5.1, III.6.10.3--6.10.4]{CohWeibelKbook}.
If a global class is zero at every finite completion, every tame
coordinate is divisible by $m$. Lift the finitely supported vector of
divisors to $K_2(\mathbb Q)$ and subtract $m$ times that lift.
The remainder lies in $K_2(\mathbb Z)$. For odd $m$ it is divisible
by $m$; for even $m$, its nonzero class is detected modulo two at
$\mathbb Q_2$, since the quaternion symbol $(-1,-1)_2$ is nontrivial.
For the latter assertion, a putative solution of $X^2+Y^2=-1$ gives,
after clearing denominators, a primitive sum of three squares equal to
zero modulo four, which is impossible. This proves injectivity of
finite-prime localization in finite weight-two cohomology. Each local
cup is zero by the Steinberg relation for $(x_q,1-x_q)$, so the global
finite cup is zero.

We retain the possible derived-limit obstruction. Set
$S=\{\ell,\infty\}$, $T=\mathbb Z_\ell(2)$, and $T_n=T/\ell^n$.
The residue map
\[
 H^1(G,T_n)\longrightarrow
       \bigoplus_{q\ne\ell}H^0(\mathbb F_q,T_n(-1))
\]
is surjective. By the purity sequence of A.22 its cokernel is the
kernel of $H^2(G_{\mathbb Q,S},T_n)\to H^2(G,T_n)$.
For odd $\ell$ the former group is zero, by the integral comparison
with $K_2(\mathbb Z)\{\ell\}=0$ and the vanishing of integral
degree three. For $\ell=2$, the integral second group is
$\mathbb Z/2$ and is detected isomorphically at the real place.
Indeed $\epsilon(g)=(\chi(g)-1)/2$ represents $\kappa(-1)$, and
$\epsilon\cup\epsilon$ at two complex conjugations has value $-1$.
It is the nonzero real class. Integral degree three is zero: in finite
coefficients the real third group is $T_n[2]$, and its generator at
level $n+1$ reduces to zero at level $n$; the finite degree-two groups
have zero derived limit. The coefficient exact sequence shows that
$H^2(G_{\mathbb Q,S},T_n)$ is also detected isomorphically at the
real place. Its map to absolute cohomology is therefore injective.
Here $K_2(\mathbb Z[1/\ell])\{\ell\}=K_2(\mathbb Z)\{\ell\}$,
since the possible new residue group $K_1(\mathbb F_\ell)$ has
order $\ell-1$; this is the same localization argument as in A.22.
These calculations use the actual two-primary comparison
\cite[VI.9.11, case $n=2$]{CohWeibelKbook} and real restriction
\cite[I.4.10(c)]{CohMilneADT}; no odd-primary theorem is applied at two.

We next prescribe integral local primitives of $a\cup b$ at
$q\ne\ell$. Write $m_q=v_q(A)$, $n_q=v_q(B)$; ordinarity makes
these actual integers in the three ordinary valuation patterns.
Principal $q$-units are uniquely $\ell$-divisible.
If $m_q>0$, then $b=\delta u$ locally; use
$h_q(g)=-a(g)\chi(g)u$.
If $n_q>0$, then $a=\delta u$ locally; use $h_q(g)=u b(g)$.
If $m_q=n_q<0$, then locally $b=a+\epsilon+\delta u$, and use
\[
 h_q(g)=-\frac{a(g)^2+\chi(g)a(g)}2-a(g)\chi(g)u.
\]
This is integral even at $\ell=2$, since $\chi(g)$ is odd.
The cocycle identities give, respectively,
\[
 \delta(-a\chi u)=a\cup\delta u,\quad
 \delta(ub)=\delta u\cup b,\quad
 \delta[-(a^2+\chi a)/2]=a\cup(a+\epsilon).
\]
These are checked symbolically in
\path{code/check_cohomology_audit.py}. When $m_q=n_q=0$, the literal
Kummer cocycles are unramified, and the procyclic unramified quotient
has zero integral second cohomology; choose an inflated primitive.

On inertia, the first two primitives take values in
$m_q\mathbb Z_\ell$ or $n_q\mathbb Z_\ell$; the negative-valuation
primitive takes values in $\ell^{-v_\ell(2)}m_q\mathbb Z_\ell$.
In the unramified case it is zero. The valuation vectors of a completed
global multiplicative class tend to zero $\ell$-adically, because
each finite power quotient has finite valuation support. Consequently
the prescribed local inertia values are zero modulo $\ell^n$ outside
a finite set of primes.

Choose any global primitive $f_n$ of the finite cup. The differences
$(f_n-h_q)|_{G_q}$ are local one-cocycles, and their inertia classes
have finite support. The residue surjection supplies a global finite
one-cocycle with those inertia classes. Subtract it from $f_n$.
Because inertia acts trivially on $T_n$ for $q\ne\ell$, equality
of those residue classes is equality of actual inertia homomorphisms.
Thus the set $P_n$ of global primitives with exactly the prescribed
inertia restrictions is nonempty. Its differences are precisely
$Z^1(G_{\mathbb Q,S},T_n)$. Indeed a difference is a global cocycle
which vanishes on every inertia group away from $\ell$. The coefficient
action is trivial on these groups, so the cocycle identity makes it
vanish on their conjugates and on their closed normal closure. It
therefore descends literally to $G_{\mathbb Q,S}$. Conversely every
such descended cocycle preserves the prescribed restrictions. This
cocycle set is finite by fixed-ramification finiteness. Reduction
preserves the specified restrictions. The
nonempty finite spaces $P_n$ therefore have nonempty inverse limit:
each finite family of compatibility equations is solved by projection
from its highest level. Their limit is a continuous integral primitive.
Taking the product over $\ell$ gives the asserted continuous
$\widehat{\mathbb Z}(2)$-primitive.

In the actual tangential coordinates, put $C=[x,y]$ and write an
element of the class-two geometric quotient as $y^b x^a C^h$.
The multiplication law is
$(a,b,h)(a',b',h')=(a+a',b+b',h+h'+ab')$; the action is
$(a,b,h)\mapsto(\chi a,\chi b,\chi^2h)$.
The latter follows from the actual path action
\cite[(12.11)--(12.15)]{CohWickelgren2012}, since its path cocycle
lies in the second commutator group. If $\delta f=a\cup b$, then
$g\mapsto y^{b(g)}x^{a(g)}C^{-f(g)}s_0(g)$ is a continuous
group-theoretic section. This proves existence in the arithmetic
quotient itself, rather than only for an abstract central extension.
\end{proof}

\begin{result}{A.26}{Proved}{Integral weight-two first cohomology}
\label{coh:h1}
There are isomorphisms
\[
 H^1_{\rm cts}(G_{\mathbb Q},\mathbb Z_\ell(2))\simeq
 \begin{cases}\mathbb Z/8,&\ell=2,\\
               \mathbb Z/3,&\ell=3,\\
               0,&\ell\geq5,
 \end{cases}
 \qquad
 H^1_{\rm cts}(G_{\mathbb Q},\widehat{\mathbb Z}(2))
       \simeq\mathbb Z/24.
\]
The last group is generated by
$\eta(g)=(\chi(g)^2-1)/24$. This confirms
GT\_TOWER\_CHECK C.1.18.
\end{result}
\begin{proof}
The finite-coefficient universal coefficient sequence is
\[
 0\to K_3(\mathbb Q)/\ell^n\to
 K_3(\mathbb Q;\mathbb Z/\ell^n)\to K_2(\mathbb Q)[\ell^n]\to0.
\]
The transition maps are reduction in the first term and multiplication
by $\ell$ in the last: the map of coefficient cofibres uses
multiplication by $\ell$ on its source and the identity on its
target. This is the natural sequence of
\cite[IV.2.5]{CohWeibelKbook}. The tame-symbol sequence shows that
$T_\ell K_2(\mathbb Q)=0$: every compatible torsion sequence has
zero image in every finite residue group, and therefore lies in the
kernel $\mathbb Z/2$, where compatibility kills it as well.
As $K_3(\mathbb Q)=\mathbb Z/48$, the finite left groups have
zero derived limit; hence
$\varprojlim K_3(\mathbb Q;\mathbb Z/\ell^n)=
 K_3(\mathbb Q)\otimes\mathbb Z_\ell$.

The natural motivic edge sequence has exact tail
\[
 K_3^M(\mathbb Q)/\ell^n\longrightarrow
 K_3(\mathbb Q;\mathbb Z/\ell^n)\longrightarrow
 H^1(G,\mathbb Z/\ell^n(2))\longrightarrow0.
\]
The source is $K_3^M(\mathbb Q)=\mathbb Z/2$, and its integral
map into $K_3(\mathbb Q)=\mathbb Z/48$ is injective
\cite[VI.4.3.1--4.3.2; VI.5.2.1]{CohWeibelKbook}.
At $\ell=2$ its image survives for $n\geq4$; at odd primes it is
zero. The finite kernels again have zero derived limit. Taking inverse
limits yields $(\mathbb Z/24)\otimes\mathbb Z_\ell$. The only
continuous-cohomology comparison error in degree one is
$\varprojlim{}^1H^0(G,\mathbb Z/\ell^n(2))$, which vanishes
because its terms are finite.

An integral cocycle into $\mathbb Q_\ell(2)$ has bounded image;
after multiplying by a power of $\ell$ it is integral. The computed
finite integral first cohomology therefore gives
$H^1(G,\mathbb Q_\ell(2))=0$. Also
$H^0(G,\mathbb Q_\ell(2))=0$. The connecting map from
$(\mathbb Q_\ell/\mathbb Z_\ell)(2)^G$ consequently identifies
this invariant subgroup with integral first cohomology. Surjectivity
of the cyclotomic character says the invariants are killed by the
ideal generated by $u^2-1$ for $u\in\mathbb Z_\ell^\times$.
That ideal is $8\mathbb Z_2$, $3\mathbb Z_3$, or
$\mathbb Z_\ell$ for $\ell\geq5$: use $u=3$ at two, $u=2$
at three, and a residue different from $\pm1$ otherwise. The class
of $1/24$ generates the invariants in each component; its connecting
cocycle is $\eta$. Continuous cochain complexes with product
coefficients are the products of the component complexes, and products
are exact in abelian groups. This proves the full assertion.
\end{proof}

\begin{result}{A.27}{Proved}{Exactly twenty-four unrestricted class-two classes}
\label{coh:24}
The fibre of geometric-conjugacy classes of continuous integral
class-two sections above any ordinary completed pair is a simply
transitive $\mathbb Z/24$-set. For the geometrically pro-$\ell$
quotient the fibre has $8$, $3$, or $1$ element according as
$\ell=2$, $\ell=3$, or $\ell\geq5$. This confirms
GT\_TOWER\_CHECK C.1.19, with ``unrestricted'' understood.
\end{result}
\begin{proof}
Existence is A.25. Fix literal abelian cocycles and one central
coordinate $h$. Every other lift is obtained by adding an integral
weight-two cocycle. Central geometric conjugation changes it by a
coboundary, giving the torsor of A.26. There are no further
identifications: a geometric conjugator between two normalized
sections has abelianization fixed by $G$, but
$H^0(G,\widehat{\mathbb Z}(1)^2)=0$. It is therefore central.
All representatives are $h+k\eta$, with $k$ modulo $24$.
\end{proof}

\begin{result}{A.28}{Proved}{What the obstruction class does and does not determine}
\label{coh:obstruction}
For a specified continuous pro-$\ell$ class-$r$ section $s$, the
obstruction to lifting it to class $r+1$ is a well-defined class
\[
 o_{r+1}(s)\in
 H^2_{\rm cts}(G_{\mathbb Q},\mathbb Z_\ell(r+1)^{b_{r+1}}).
\]
It vanishes exactly when that specified section lifts. Its finite
absolute reductions determine it only modulo
$\varprojlim{}^1H^1(G,\mathbb Z/\ell^n(r+1)^{b_{r+1}})$.
In particular, no calculation of a nonzero element of the ambient
derived-limit kernel, such as A.23 or A.24, proves nonvanishing of
$o_{r+1}(s)$ for an ordinary ghost.
\end{result}
\begin{proof}
Lift the geometric cocycle of $s$ through the continuous Hall set
section. Its defect is a continuous normalized two-cocycle in the
central kernel $L_{r+1,\ell}$. The action on that kernel is the
original graded $G$-action because geometric inner conjugation acts
trivially there. Changing the set lift adds a continuous coboundary.
If the defect is a coboundary, correcting by its continuous primitive
gives a section; conversely a genuine lift supplies such a primitive.
This proves the first assertion. The second is the coefficient-limit
exact sequence used in A.24, applied componentwise. The last assertion
is the precise distinction between the existence of a class in a group
and identifying the class of a specified extension. It supplies no
nonvanishing theorem for that specified extension.
\end{proof}

\begin{result}{A.29}{Proved}{Two local comparisons detect all twenty-four twists}
\label{coh:local24}
For two global integral class-two section classes with the same
abelian completed pair, equality of their geometrically pro-two
restrictions at $q=3$ and of their geometrically pro-three
restrictions at $q=2$ implies equality of the global classes.
Thus at most one of the twenty-four unrestricted global classes can
match a fixed family of local point-section classes.
\end{result}
\begin{proof}
After literal abelian normalization the difference is a class in
$H^1(G,\widehat{\mathbb Z}(2))$. At $q\ne\ell$, integral
weight-two cocycles have zero inertia by the relation
$(q^2-q)c(T_q)=0$. Hence
\[
 H^1(G_{\mathbb Q_q},\mathbb Z_\ell(2))
           =\mathbb Z_\ell/(q^2-1)\mathbb Z_\ell,
\]
where an unramified cocycle is evaluated at arithmetic Frobenius.
At $(q,\ell)=(3,2)$ the value
$\eta(F_q)=(9-1)/24=1/3$ generates $\mathbb Z_2/8$.
At $(q,\ell)=(2,3)$ the value $(4-1)/24=1/8$ generates
$\mathbb Z_3/3$. Therefore the two local maps together are
injective on the global $\mathbb Z/24$.
Moreover $H^0(G_{\mathbb Q_q},\mathbb Z_\ell(1)^2)=0$, since
$q-1$ is a nonzero scalar on a torsion-free module. Thus local
geometric conjugacy between two normalized sections is already central,
and the local fibre action by this first cohomology group is free.
The injectivity calculation therefore applies to section classes,
not merely to chosen central coordinates.
\end{proof}

\subsection{Local matching at a nonzero ordinary valuation}

\begin{result}{A.30}{Proved}{Rational formal conjugacy detects pro-$\ell$ conjugacy}
\label{coh:magnusconj}
Embed the free pro-$\ell$ group $F_\ell=\langle x,y\rangle^{(\ell)}$
by $x\mapsto1+X$, $y\mapsto1+Y$ in the noncommutative formal
power-series ring over $\mathbb Z_\ell$. If two group elements
are conjugate by a unit of
$\mathbb Q_\ell\langle\!\langle X,Y\rangle\!\rangle$, they
are conjugate in $F_\ell$. This confirms GT\_TOWER\_CHECK C.1.16.
\end{result}
\begin{proof}
At each homogeneous degree, sum coefficients over every orbit of words
under cyclic rotation. This map has equal values on $AB$ and $BA$:
each concatenation of two words is cyclically equivalent to the reversed
concatenation. It consequently has equal values on $UgU^{-1}$ and
$g$, for every invertible formal series $U$. If the Magnus series of
$g,h\in F_\ell$ are rationally conjugate, every cyclic coefficient
sum of $g-h$ is zero. Within each finite orbit of words, an integral
vector of coefficient sum zero is an integral linear combination of
differences of words in that orbit. A difference between two cyclic
rotations is a sum of additive commutators $AB-BA$. Thus every finite
homogeneous truncation of $g-h$ lies in the integral additive
commutator module.

Let $P$ be any finite $\ell$-group quotient of $F_\ell$.
The augmentation ideal of $\mathbb F_\ell[P]$ is nilpotent. To see
this directly, choose a central element $z$ of order $\ell$, apply
induction to $P/\langle z\rangle$, and use $(z-1)^\ell=0$;
the base case is the trivial group. Therefore a power of the integral
augmentation ideal lies in $\ell\mathbb Z_\ell[P]$, and its
powers tend to zero $\ell$-adically. Evaluation
$X\mapsto\bar x-1$, $Y\mapsto\bar y-1$ consequently gives a
continuous ring homomorphism from integral formal series to
$\mathbb Z_\ell[P]$.

The quotient of $\mathbb Z_\ell[P]$ by its additive commutator
module is free on the conjugacy classes of $P$. Indeed $ab-ba$ is
the difference of two conjugate group elements, and a difference
between an element and one of its conjugates is an additive
commutator. The module is closed, since this quotient is a finite
free $\mathbb Z_\ell$-module. Evaluation of the truncations shows
that $\bar g$ and $\bar h$ have the same conjugacy-class coordinate,
and are thus conjugate in $P$.

For each open normal $N\triangleleft F_\ell$ the set
$\{v:v^{-1}gvh^{-1}\in N\}$ is nonempty and closed.
Finite intersections contain the set for the intersection of the
corresponding $N$. Compactness of $F_\ell$ supplies an actual
conjugator. We only evaluated integral formal series in finite group
algebras; the rational conjugator was never subjected to that
evaluation. The Magnus embedding is injective: if a group element
had series $1$, its evaluation in every finite $\ell$-group algebra
would be $1$, so its image in every finite group quotient would be
the identity.
\end{proof}

\begin{result}{A.31}{Proved}{The local boundary fibre is a singleton}
\label{coh:boundarylocal}
Let $q\ne\ell$ be rational primes and let
$t\in\mathbb Q_q\setminus\{0,1\}$ have
$(v_q(t),v_q(1-t))\ne(0,0)$. In the actual geometrically
pro-$\ell$ arithmetic fundamental group of the tripod over
$\mathbb Q_q$, the geometric-conjugacy fibre above
$(\kappa_\ell(t),\kappa_\ell(1-t))$ consists only of the
point-section class of $t$. Therefore a global full pronilpotent
section with ordinary abelian pair matches its prescribed local
point at every support prime $q$ in every coefficient component
$\ell\ne q$. This confirms GT\_TOWER\_CHECK C.1.17. It says
nothing about the equal residue/coefficient prime or zero-valuation
places.
\end{result}
\begin{proof}
The actual arithmetic action on $F_\ell$ is unramified at $q$.
One can check this without replacing the action by a cyclotomic
model. At class two the formula in A.25 is explicitly unramified.
Inductively the kernel of the automorphism map from class $i+1$ to
class $i$ embeds in
$\operatorname{Hom}(N_{1,\ell},L_{i+1,\ell})$, on which
Frobenius has weight $q^i$. An inertia cocycle there is killed by
$q^i-q$, using the tame relation; it vanishes for $i\geq2$ because
the module is torsion free. Wild and the other tame factors have no
homomorphisms to this additive pro-$\ell$ module. The induction,
followed by A.21, proves the assertion about the full action.

Suppose first $m=v_q(t)>0$. The second local Kummer class is zero,
as $1-t$ is a principal unit. Every section cocycle factors through
the tame quotient
$\mathbb Z_\ell(1)\rtimes\widehat{\mathbb Z}$: on the
discarded inertia kernel the action is trivial and a homomorphism
from the pro-prime-to-$\ell$ kernel into $F_\ell$ is zero.
Write $V$ for its geometric tame inertia image and $z$ for its
geometric Frobenius coordinate. If $\sigma_0$ denotes the actual
tangential Frobenius automorphism, then
\[
      z\sigma_0(V)z^{-1}=V^q,\qquad
      V_{\rm ab}=(m,0),\qquad \sigma_0(x)=x^q.
\]
The nonzero ordinary integer $m$ is nonzero in $\mathbb Q_\ell$.

Pass to the rational prounipotent group obtained from the Hall
polynomials, retaining the actual automorphism $\sigma_0$.
There is a unique $u$ with $u^{-1}z\sigma_0(u)=1$.
To prove this, solve in successive central weights: the correction
in weight $i$ is multiplication by $q^i-1$, an invertible scalar
over $\mathbb Q_\ell$. The compatible corrections give $u$.
For $W=u^{-1}Vu$ the equation becomes $\sigma_0(W)=W^q$.
The element $x^m$ satisfies it with the same weight-one coordinates.
If their logarithms first differ in weight $i\geq2$, that difference
is killed by $q^i-q$, another invertible rational scalar. Induction
shows $W=x^m$. Thus $V$ is conjugate to $x^m$ by a rational formal
unit, and A.30 supplies a conjugator in $F_\ell$ itself.

After this integral conjugation the Frobenius coordinate $z'$
centralizes $x^{qm}$. Its centralizer in $F_\ell$ is
$x^{\mathbb Z_\ell}$. Here is a direct verification. Over rational
formal series, $X$ is a formal series in $(1+X)^{qm}-1$, since the
linear coefficient $qm$ is nonzero. An element centralizing
$x^{qm}$ therefore centralizes $X$. Comparing coefficients in
$XA=AX$ moves terminal $X$'s to the front; a word ending in $Y$
is then forced to have coefficient zero. Every word containing $Y$
has coefficient zero, so $A\in\mathbb Q_\ell[[X]]$. The group
retraction $y\mapsto1$ and Magnus injectivity show that an element
of $F_\ell$ with this property lies in $x^{\mathbb Z_\ell}$.

Consequently the section cocycle lies in that procyclic subgroup,
as its inertia and Frobenius do. The equivariant inclusion
$\mathbb Z_\ell(1)\to F_\ell$, $a\mapsto x^a$, has the first
abelian coordinate as a retraction. Two such cyclic cocycles with
the same Kummer class differ by a cyclic coboundary and hence by
geometric conjugation. Applying the argument to the actual point
section of $t$ proves the singleton assertion.

If $v_q(1-t)>0$, use the curve automorphism $t\mapsto1-t$.
If both valuations are negative, use $t\mapsto1/t$, whose abelian
map is $(A,B)\mapsto(A^{-1},-BA^{-1})$ and whose first valuation
is positive. These curve automorphisms induce isomorphisms of the
actual arithmetic fundamental groups; changing the transported base
path changes them only by geometric inner automorphism. Thus they
give bijections of the relevant section-class fibres. Finally, a
global pronilpotent section projects continuously to each pro-$\ell$
quotient and restricts continuously to $G_{\mathbb Q_q}$, so the
local assertion applies at each stated pair of primes.
\end{proof}

\paragraph{An independent published local comparison.}
\citet[Proposition~A.6, author pp.~37--38]{RootHoshi2014Conditional}
places ramified inertia in a cusp subgroup at good reduction away
from the coefficient prime. Its peripheral-recognition input is
\citet[Lemma~1.6, author pp.~15--16]{RootHoshiMochizuki2011}.
After that recognition, the cyclic-centralizer and Kummer arguments
above give the same singleton conclusion. Both primary proofs were
opened at those locators. This corroborates A.31; it supplies no
additional restriction on unramified section classes at zero-support
places.

